% !TeX root = GeoDynamics.tex
\section{Geo GLSL Export Process}

\subsection{Purpose}

The Geo GLSL export process exists to keep the Python dynamics model and the
Vulkan shader implementation aligned without hand-maintaining two independent
versions of the same collision logic.

The export target for dynamics is:

\begin{verbatim}
C:\_DJ\gPCD\vulkan\sim\common\GeoDynamics.glsl
\end{verbatim}

\texttt{GeoDynamics.glsl} is a generated file.  It should be regenerated from
the current Geo source model whenever the exported dynamics changes.  It should
not be manually patched as the normal workflow.

\subsection{Source Of Truth}

The source of truth for Geo collision dynamics is the Python Geo path:

\begin{itemize}
    \item \texttt{base/GeoDynamics.py} owns GLSL-portable dynamics functions.
    \item \texttt{base/GeoBase.py} owns the Python harness, configuration
    loading, particle creation, direct all-pairs scan, and frame loop.
    \item \texttt{base/Reporting.py} owns reporting and diagnostics.
\end{itemize}

The exported shader should preserve the shape of the GLSL-portable functions in
\texttt{GeoDynamics.py}.  Python-only harness behavior must not be exported into
\texttt{GeoDynamics.glsl}.

\subsection{Generated And Hand-Owned Files}

The following file is generated by the Geo export process:

\begin{itemize}
    \item \verb|C:\_DJ\gPCD\vulkan\sim\common\GeoDynamics.glsl|
\end{itemize}

The following files are hand-owned integration files.  The exporter may check
them, but it should not silently rewrite them unless that behavior is explicit:

\begin{itemize}
    \item \verb|C:\_DJ\gPCD\vulkan\sim\common\particle.glsl|
    \item \verb|C:\_DJ\gPCD\vulkan\src_app\particleOnly\ParticleStruct.hpp|
    \item \verb|C:\_DJ\gPCD\vulkan\sim\common\move_particle.glsl|, unless
    motion is moved into the generated dynamics file.
    \item \verb|C:\_DJ\gPCD\vulkan\sim\BoundarySpheresMotion\BoundarySpheresMotionP.comp|
    \item \verb|C:\_DJ\gPCD\vulkan\sim\BoundarySpheresMotion.lst|
\end{itemize}

\subsection{Export Contract}

Each export must produce a complete \texttt{GeoDynamics.glsl}.  The generated
file must be usable after deleting the previous generated file.  The exporter
must not depend on local manual edits that survive from an earlier export.

The generated file should define the shader-side equivalents of the Geo
dynamics entry points used by the compute shader:

\begin{verbatim}
GeoBeginContactFrame(SourceID)
GeoAddParticleContact(SourceID, TargetID)
GeoProcessCollisions(SourceID)
GeoMoveParticle(SourceID, positionBuffer, dt)
\end{verbatim}

If contact clearing remains in the compute shader instead of dynamics, the
exporter must document that choice and the compute shader must not call a
missing \texttt{GeoBeginContactFrame} function.

\subsection{Particle Layout Contract}

The host-side particle layout and shader-side particle layout must stay aligned.
The neutral contact names are:

\begin{verbatim}
ContactState contacts[MAX_CONTACTS]
contactCount
\end{verbatim}

The C++ definition in \texttt{ParticleStruct.hpp} and the GLSL definition in
\texttt{particle.glsl} must use the same field order and compatible vector
types.  Field names should be generic when they describe shared storage.  Names
such as \texttt{NeoContactState}, \texttt{GeoContactState}, \texttt{ncs}, and
\texttt{gcs} should not be reintroduced for the shared particle buffer.

\subsection{Export Steps}

The normal export workflow is:

\begin{enumerate}
    \item Update \texttt{GeoDynamics.py} and, when needed, \texttt{GeoBase.py}.
    \item Run the Geo Python checks for the target case, such as
    \texttt{TwoParticleHorizontal.cfg}.
    \item Regenerate \texttt{GeoDynamics.glsl} from the exporter.
    \item Confirm \texttt{BoundarySpheresMotionP.comp} includes
    \texttt{../common/GeoDynamics.glsl}, not obsolete dynamics files.
    \item Confirm \texttt{BoundarySpheresMotion.lst} lists
    \texttt{common/GeoDynamics.glsl}.
    \item Compile the compute shader with \texttt{glslc}.
    \item Run the Vulkan target configuration.
\end{enumerate}

\subsection{Verification}

After export, the first verification is a shader compile:

\begin{verbatim}
C:\_DJ\gPCD3rdParty\VulkanSDK\Bin\glslc.exe ^
  C:\_DJ\gPCD\vulkan\sim\BoundarySpheresMotion\BoundarySpheresMotionP.comp ^
  -o %TEMP%\BoundarySpheresMotionP.spv
\end{verbatim}

The compile must not reference deleted or obsolete files such as
\texttt{NeoDynamics.glsl}.  It must also not fail on missing Geo entry points.

The second verification is a Vulkan run using the exported case configuration,
for example:

\begin{verbatim}
C:\_DJ\gPCD\python\cfg_gendata\TwoParticleHorizontal.cfg
\end{verbatim}

\subsection{Do Not Patch Generated Dynamics}

When \texttt{GeoDynamics.glsl} is wrong, fix the exporter or the Python source
model and regenerate the file.  Manual edits to \texttt{GeoDynamics.glsl} are
allowed only as short-lived experiments.  Once the experiment works, the change
must be moved back into the export source and the file must be regenerated.

This rule keeps the generated shader disposable.  A valid export process should
be able to delete \texttt{GeoDynamics.glsl}, regenerate it, compile the shader,
and run the target case without recovering hidden hand edits.
